%% kapa: exam=1 round=25 year=2014 period=1 form=A subject=economics subject_ko=경제원론 seq=15 number=55
%% source: Q-Net 감정평가사 기출문제 (cst00309) artlSeq=5101207 file=extracted/1008/5101207/2014년 제25회 감정평가사 1차시험 기출문제/1교시(A형).hwp
%% conversion: kordoc 4.24.51 -> tools/gen1.py (auto); review=proofread-sonnet; flags=none
\begin{question}{55}
민간항공기는 A국과 B국에서만 생산한다. 두 국가 항공사의 보수(이윤)행렬이 다음과 같다. 만일 두 나라 정부는 자국 항공사가 생산할 경우 각각 10억 달러의 생산보조금을 지급할 때 내쉬균형은 무엇인가? \par
{\small\centering
\begin{tabular}{|c|c|c|c|}
\hline
\multirow{4}{*}{\makecell{A국\newline{}항공사}} & \multirow{2}{*}{} & \multicolumn{2}{|c|}{\makecell{B국 항공사}} \\
\cline{3-4}
 &  & \makecell{생산함} & \makecell{생산하지 않음} \\
\cline{2-4}
 & \makecell{생산함} & \makecell{A국 -5억 달러\newline{}B국 -5억 달러} & \makecell{A국 100억 달러\newline{}B국 0 달러} \\
\cline{2-4}
 & \makecell{생산하지 않음} & \makecell{A국 0 달러\newline{}B국 100억 달러} & \makecell{A국 0 달러\newline{}B국 0 달러} \\
\hline
\end{tabular}\par}

\begin{choices}
  \item A국과 B국에서 생산함
  \item A국에서만 생산함
  \item B국에서만 생산함
  \item A국과 B국에서 생산하지 않음
  \item ②, ④ 두 가지 경우
\end{choices}
\end{question}
\answer{1}
